\chapter{Scaling Laws and Compute Budgets}
\label{ch:scaling-laws}

% Part II, Chapter 8: Choosing model size, token count, and training duration.
%
% Core question: Given a fixed GPU-hour budget, how should you allocate it
% between model size and training data?
%
% Planned sections:
%   8.1  Hook: "You Have \$10M of GPU Time. How Do You Spend It?"
%   8.2  Power Laws: Loss as a Function of N, D, and C
%   8.3  Kaplan et al. (2020): Bigger Models Are More Sample-Efficient
%   8.4  Chinchilla (Hoffmann 2022): Compute-Optimal Training
%   8.5  Overtraining: Why LLaMA Trained 4x Beyond Chinchilla-Optimal
%   8.6  Inference Cost vs Training Cost: The Deployment Tradeoff
%   8.7  Emergent Capabilities: What Appears at Scale
%   8.8  Pilot Runs and Scaling Predictions
%   8.9  Failure Modes: Misreading Scaling Curves
%
% Key thesis: Scaling laws are not just empirical curiosities. They are
% resource-allocation tools. The difference between a good and bad scaling
% decision can be millions of dollars of wasted compute.
%
% Lab 8 (planned): Fit scaling laws from pilot runs at 3--4 model sizes,
% predict final loss for a target compute budget, then verify.

\textit{This chapter is a placeholder. It will cover Kaplan scaling laws,
Chinchilla compute-optimal training, overtraining regimes, inference vs
training cost tradeoffs, emergent capabilities, and how to use pilot runs
to make budget decisions.}
